\documentclass[11pt]{article}
\usepackage{fontspec}
\usepackage[french,english]{babel}
\usepackage{amsmath}
\usepackage{amssymb}
\usepackage{amsthm}
\usepackage{mathrsfs}
\usepackage{tikz-cd}
\usetikzlibrary{babel}
\usepackage{enumitem}

\newcommand{\printedpage}[1]{\marginpar{\footnotesize p.~#1}}
\newcommand{\msp}[1]{\marginpar{\footnotesize [#1]}}
\newcommand{\Ind}{\operatorname{Ind}}
\newcommand{\uHom}{\underline{\mathrm{Hom}}}
\newcommand{\Ob}{\operatorname{Ob}}
\newcommand{\card}{\operatorname{card}}
\newcommand{\Fl}{\mathrm{Fl}}
\newcommand{\PF}{\mathrm{PF}}

\begin{document}

\printedpage{24}
% head: A. Grothendieck : Dérivateurs, Ch. XVIII, Édition des Universités de Montpellier II et Paris VII
% foot: Version 17 décembre 2025

\textit{with}
\[
\Psi(i) \;=\; \varinjlim_{j\in J}{}_{C}\,\psi(i,j)\;.
\]
\textit{c)} $\Ind_\pi C$ \textit{is stable under small direct limits} (\footnote{also add, in d), corollary 4.4.16, page 101.}).

It remains to establish this last point. Since we already have finite $\varinjlim$ by b, it remains to prove the existence of small filtering $\varinjlim$, or, if preferred, of infinite sums, which will be nicer in fact. But
\msp{page 97}
we can just as well treat the general case of a $\varinjlim$ of arbitrary type $J$ ($J$ small). For this, let
\[
I \;=\; I(J) \;=\; \begin{aligned}[t]
&\text{set of subcategories } J' \text{ of } J \text{ such that}\\
&\card\Fl(J')\leqslant\pi,
\end{aligned}
\]
let us order $I$ by inclusion. It is obvious that $I\neq\emptyset$ (for $\emptyset\in I$), and even that $I$ \textit{is large relative to} $\pi$, for if $(J'_\alpha)_{\alpha\in A}$ with $\card(A)\leqslant\pi$ is a family of elements of $I$, then the subcategory of $J$ generated by the $J'_\alpha$ is still in $I$ (immediate verification), and is an upper bound for the $J'_\alpha$. In fact, the consideration of $I(J)$ associated with $J$ will allow us to prove this, which generalizes 4.4.13 c), taking b) into account:

\textbf{Corollary 4.4.14} (\footnote{Compare also 4.11.3, page 270.}). \textit{Let $\mathcal{M}$ be a category, $\pi$ an infinite cardinal such that}
\begin{enumerate}
\item[a)] \textit{$\mathcal{M}$ is}
\msp{page 98}
\textit{stable under $\varinjlim$ of cardinality $\leqslant\pi$, and}
\item[b)] \textit{$\mathcal{M}$ is stable under $\varinjlim$ over ordered sets $I$ large relative to $\pi$.}
\end{enumerate}
\textit{Then $\mathcal{M}$ is stable under small $\varinjlim$ .}

(\footnote{\textbf{Corollary 4.4.14.1.} $\mathcal{M}$ being as in 4.4.14, and $f:\mathcal{M}\longrightarrow\mathcal{E}$ being a functor which commutes with $\varinjlim$ of cardinality $\leqslant\pi$, and with small $\varinjlim$ large relative to $\pi$, then $f$ commutes with small $\varinjlim$ .}).

Indeed, if $\varphi:J\longrightarrow\mathcal{M}$ is given, we construct from it
\[
\psi\;:\;I\;=\;I(J)\;\longrightarrow\;\mathcal{M}
\]
by
\[
\psi(J')\;=\;\varinjlim_{J'}\varphi(J')\;,
\]
which exists by the hypothesis on $\mathcal{M}$. It is then immediate that $\varinjlim\limits_{J}\varphi$ exists if and only if $\varinjlim\limits_{I}\psi$ exists, now this existence is assured by the fact that $I$ is large relative to $\pi$, and the hypothesis on $\mathcal{M}$.

To finish proving theorem 4.4.13, we are thus reduced to the particular case of direct limits $\varinjlim\limits_{J}\varphi(j)$, when $\varphi:J\longrightarrow\Ind'_\pi C$ is given, with $J$ an
\printedpage{25}
% head: A. Grothendieck : Dérivateurs, Ch. XVIII, Édition des Universités de Montpellier II et Paris VII
% foot: Version 17 décembre 2025
\msp{page 99}
ordered set large relative to $\pi$. But this will follow from the more precise result:

\textbf{Corollary 4.4.15} (\footnote{compare with 4.7.3.1, which ought to be made explicit here.}). \textit{Let $C$ be a small category. Then the strictly full subcategory} $\Ind'_\pi C$ \textit{of $C^\wedge$ is stable in $C^\wedge$ under $\varinjlim$ large relative to $\pi$, in other words, these $\varinjlim$} exist in $\Ind'_\pi(C)$, \textit{and the inclusion functor commutes with them.} % doubt resolved: existent is upright in the PDF; see corrections.csv, PDF leaf 27

\textsc{Proof}. Let $I$ be an ordered set, large relative to $\pi$, and $i\longmapsto F_i$ a functor $I\longrightarrow\Ind'_\pi(C)$, $F\;=\;\varinjlim\limits_{I}F_i$ in $C^\wedge$, let us prove that $F\in\Ob\Ind'_\pi(C)$, i.e. (4.4.3) that $C/F$ is large relative to $\pi$. I use criterion b) of 4.4.4, we have to check $\PF_\pi\,1$, $\PF_\pi\,2$. Let
\[
x_\alpha\xrightarrow{\;u_\alpha\;}F\;=\;\varinjlim F_i
\]
be objects of $C/F$ indexed by $A$ with $\card A\leqslant\pi$. Each $u_\alpha$ comes from $x_\alpha\longrightarrow F_{i_\alpha}$, and taking a common upper bound for the $i_\alpha$ ($I$ large relative to $\pi$), we find an $i\in I$ and $v_\alpha:x_\alpha\longrightarrow F_i$. Since $F_i$ is in $\Ind'_\pi(C)$, i.e. $C/F_i$ large relative to $\pi$, considering the $v_\alpha$ as objects of $C/F_i$, they are bounded above by an object $x\longrightarrow F_i$ of $C/F_i$. We thus have
\[
\begin{tikzcd}
x_\alpha \arrow[r] \arrow[d] & F_{i_\alpha} \arrow[d] \arrow[dr, dotted] & \\
x \arrow[r] & F_i \arrow[r] & F\;,
\end{tikzcd}
\]
which shows that the $x_\alpha\longrightarrow F$ are bounded above by $x\longrightarrow F$,
\msp{page 100}
which proves $\PF_\pi\,1$. Let us prove $\PF_\pi\,2$, i.e. let $x\xrightarrow{\;u\;}F$, $y\xrightarrow{\;v\;}F$ be two objects of $C/F$, and $(f_\alpha)_{\alpha\in A}$ a family of cardinality $\leqslant\pi$ of arrows from $u$ to $v$. We know that $u$, $v$ come from $u_i:x\longrightarrow F_i$, $v_i:y\longrightarrow F_i$ for $i$ large enough.
\[
\begin{tikzcd}[column sep=small]
x \arrow[rr, "f_\alpha"] \arrow[dr, "u"'] & & y \arrow[dl, "v"] \\
& F &
\end{tikzcd}
\]
\[
\begin{tikzcd}[column sep=small]
x \arrow[rr, "f_\alpha"] \arrow[dr, "u_i"'] & & y \arrow[dl, "v_i"] \\
& F_i &
\end{tikzcd}
\]

For every $\alpha$, we can choose $i_\alpha\geqslant i$ such that $v_if_\alpha$ and $u_i$ become equal on composing with $F_i\longrightarrow F_{i_\alpha}$, and since $\card A\leqslant\pi$ and $I$ large relative to $\pi$, there exists an upper bound $i'$ for the $i_\alpha$. All told, we can assume (replacing $i$ by $i'$) that $v_if_\alpha=u_i$ for every $\alpha$. Then $u_i$, $v_i$ being considered as objects of $C/F_i$, the $f_\alpha$ are arrows of $C/F_i$ from $u_i$ to $v_i$. Since $C/F_i$ is large relative to $\pi$ (by the hypothesis $F_i\in\Ob\Ind'_\pi(C)$), we find that there exists an equalizer $(y,v_i)\xrightarrow{\;g\;}(z,w_i)$.
\printedpage{26}
% head: A. Grothendieck : Dérivateurs, Ch. XVIII, Édition des Universités de Montpellier II et Paris VII
% foot: Version 17 décembre 2025
\[
\begin{tikzcd}
x \arrow[r, shift left=2, "(f_\alpha)"] \arrow[r] \arrow[r, shift right=2] \arrow[dr, "u_i"'] & y \arrow[r, "g"] \arrow[d, "v_i"'] & z \arrow[dl, "w_i"] \\
& F_i &
\end{tikzcd}
\]

Then the image of $g$ in $C/F$ (via $F_i\longrightarrow F$) equalizes the arrows $f_\alpha$, as arrows of $C/F$, which proves $\PF_\pi\,2$ and completes the proof of 4.4.14, and thereby also that of theorem 4.4.13, and of lemma 4.3. This also completes the proof of the first assertion of lemma
\msp{page 101}
4.3.7 (p. 77), which completes the proof of theorem 4.3.6 (p. 70-71). It remains to establish the universal property stated in 4.3.7, which I am going to recall:

\textbf{Corollary 4.4.16} (\footnote{Could be incorporated as part d) in theorem 4.4.13, for it is an important complement, and not so obvious after all.}). \textit{Let $C$ be a category equivalent to a small category, stable under $\varinjlim$ of cardinality $\leqslant\pi$, where $\pi$ is a given infinite cardinal, and let us consider the canonical inclusion functor}
\[
C\xrightarrow{\;\alpha\;}\Ind_\pi(C)\;. \tag*{(4.4.16.1)}
\]
\textit{This functor is} 2\textit{-universal for functors $f:C\longrightarrow\mathcal{E}$, where $\mathcal{E}$ is a $\mathcal{U}$-category stable under small $\varinjlim$ , and $f$ a functor commuting with these direct limits. In other words, we have}
\begin{enumerate}
\item[\textit{1$^\circ$)}] \textit{$\alpha$ has the preceding properties, i.e.} $\Ind_\pi(C)$ \textit{is stable under $\varinjlim$ , and the functor $\alpha$ commutes with $\varinjlim$ of cardinality $\leqslant\pi$.}
\item[\textit{2$^\circ$)}] \textit{For every $\mathfrak{U}$-category $\mathcal{E}$ stable under small direct limits, the functor}
\[
\begin{aligned}
g\;&\longmapsto\;g\circ\alpha\\
\uHom_{!}(\Ind_\pi C,\mathcal{E})\;&\longrightarrow\;\uHom_{!\,\pi}(C,\mathcal{E})
\end{aligned}
\tag*{(4.4.16.2)}
\]
\textit{is an equivalence of categories, where the subscript $!$ (resp. $!\,\pi$) indicates the strictly full subcategory of the category} $\uHom$
\msp{page 102}
\textit{under consideration, formed by the functors which commute with small $\varinjlim$ (resp. with those of cardinality $\leqslant\pi$).}
\end{enumerate}

We can thus also say that $\Ind_\pi(C)$ plays the role of an ``envelope'' stable under small $\varinjlim$ of the category $C$, when the $\varinjlim$ of cardinality $\leqslant\pi$, which already exist in $C$, are respected by the embedding of $C$ in its envelope. It is a by no means tautological result that it suffices for this to add to $C$ some $\varinjlim$ \textit{large relative to} $\pi$ of objects of $C$, and that more precisely, the envelope category is realized by $\Ind_\pi C$ (or, if preferred and less plethorically, by $\Ind'_\pi C\subset C^\wedge$). % [sic] les plongement -> le plongement, number agreement

\textsc{Proof}. 1°) is contained in 4.4.13 a) and c). Let us therefore prove 2°). It is obvious that $\alpha$ is a fully faithful functor whose essential image strictly generates
\printedpage{27}
% head: A. Grothendieck : Dérivateurs, Ch. XVIII, Édition des Universités de Montpellier II et Paris VII
% foot: Version 17 décembre 2025
$\mathcal{M}=\Ind_\pi C$, whence it follows that the functor (4.4.16.2) is fully faithful, being induced by the fully faithful functor
\[
\alpha^*_{\mathcal{E}}\;:\;\uHom(\mathcal{M},\mathcal{E})\;\longrightarrow\;\uHom(C,\mathcal{E})
\]
\msp{page 103}
on the full subcategories under consideration. It thus remains to show that (4.4.16.2) is essentially surjective. So let $f:C\longrightarrow\mathcal{E}$ be in $\uHom_{!\,\pi}(C,\mathcal{E})$, let us consider its canonical extension % [sic] les sous-catégorie pleines -> les sous-catégories pleines, number agreement
\[
\left\{
\begin{aligned}
&Lf\;:\;\Ind(C)\;\longrightarrow\;\mathcal{E}\\
&Lf((X_i)_{i\in I})\;=\;\varinjlim_{i\in I}f(X_i)
\end{aligned}
\right.
\tag*{(4.4.16.3)}
\]
(cf. SGA 4 I 8.7.3), and let
\[
\psi(f)\;=\;Lf|\Ind_\pi(C)\;. \tag*{(4.4.16.4)}
\]
Let us recall that $Lf$ commutes with small filtering $\varinjlim$, but since the inclusion functor
\[
\Ind_\pi C\;\longrightarrow\;\Ind(C) \tag*{(4.4.16.5)}
\]
does not commute with these $\varinjlim$ , it is not clear a priori that $\psi(f)$ commutes with the said limits, nor a fortiori that it commutes with small $\varinjlim$ without more. But we have seen that (4.4.16.5) commutes with $\varinjlim$ \textit{large relative to} $\pi$, which is contained in corollary 4.4.15 (since the \textit{filtering} direct limits in $\Ind(C)$ are the limits \textit{in} $C^\wedge$): to say that $\Ind'_\pi(C)\longrightarrow\Ind'(C)$ commutes with a certain type of filtering $\varinjlim$, amounts to saying that this is so for the composite with $\Ind'_\pi(C)\hookrightarrow C^\wedge$, which is established by 4.4.15. From this it follows that $\psi(f)$ \textit{commutes with} $\varinjlim$ \textit{large relative to} $\pi$. To establish that $g=\psi(f)$ commutes
\msp{page 104}
with small $\varinjlim$ , it thus remains to see that it commutes with $\varinjlim$ of cardinality $\leqslant\pi$, by virtue of 4.4.14.1 (margin page 98). To establish this, we use theorem 4.4.13 b, representing a functor $\varphi:J\longrightarrow\Ind_\pi(C)$ (where $\card\Fl J\leqslant\pi$) in indexed form by
\[
\Phi\;:\;I\times J\;\longrightarrow\;C\;, \tag*{(4.4.16.6)}
\]
$I$ being an ordered set large relative to $\pi$. We thus have
\[
\varphi(j)\;=\;\varinjlim_{i\in I}{}_{\Ind_\pi C}\,\Phi(i,j)\;,
\]
and since $g=\psi(f)$ commutes with $\varinjlim$ of type $I$ by what we have seen, we will have
\[
g(\varphi(j))\;=\;\varinjlim_{i\in I}{}_{(\mathcal{E})}\;g(\Phi(i,j))\;,
\]
whence
\[
\varinjlim_{j\in J}{}_{(\mathcal{E})}\,g(\varphi(j))\;=\;\varinjlim_{j\in J}\;\varinjlim_{i\in I}\;\underbrace{g(\Phi(i,j))}_{=\,f(\Phi(i,j))}\qquad\text{(limits in $\mathcal{E}$)}\;. \tag*{$(*)$}
\]
\printedpage{28}
% head: A. Grothendieck : Dérivateurs, Ch. XVIII, Édition des Universités de Montpellier II et Paris VII
% foot: Version 17 décembre 2025

On the other hand, we have
\[
\begin{aligned}
\varinjlim_{j\in J}{}_{\Ind_\pi}\,\varphi(j)\;&\simeq\;\varinjlim_{j\in J}\;\varinjlim_{i\in I}\;\Phi(i,j)\\
&\simeq\;\varinjlim_{i\in I}\;\underbrace{\varinjlim_{j\in J}\;\Phi(i,j)}_{\substack{\text{can be computed}\\ \text{in $C$, say}\\ \Psi(i)\,\in\,\Ob C}}\;,
\end{aligned}
\]
thus
\[
\varinjlim_{j\in J}{}_{\Ind_\pi}\,\varphi(j)\simeq\varinjlim_{i\in I}\;\Psi(i)\;, \tag*{$(**)$}
\]
\msp{page 105}
whence by applying $g$, which commutes with $\varinjlim$ of type $I$, and now recalling that $g|C\simeq f$,
\[
g(\,\varinjlim_{J}{}_{\Ind_\pi C}\,\varphi(j))\;\simeq\;\varinjlim_{i\in I}\;f(\Psi(i))\;. \tag*{$(*{*}*)$}
\]
We have on the other hand
\[
\Psi(i)\;\overset{\text{def}}{=}\;\varinjlim_{j\in J}\;\Phi(i,j)\;,\qquad\qquad \varinjlim\;\text{ in } C\;,
\]
and since $f$ commutes with $\varinjlim$ of type $J$ (for $\card\Fl J\leqslant\pi$), we conclude from this
\[
f(\Psi(i))\;\simeq\;\varinjlim_{j\in J}\;f(\Phi(i,j))\;,
\]
thus by $(*{*}*)$
\[
g(\,\varinjlim_{j\in J}\;\varphi(j))\;\simeq\;\varinjlim_{i\in I}\;\varinjlim_{j\in J}\;f(\Phi(i,j))\;. \tag*{$(*{*}{*}*)$}
\]
By composing with $(*)$ and using again the commutativity of $\varinjlim\limits_{I}$ with $\varinjlim\limits_{J}$ , we thus find
\[
g(\,\varinjlim\;\varphi(j))\;\overset{\sim}{\longleftarrow}\;\varinjlim\;g(\varphi(j))\;,
\]
which proves that $g$ commutes with direct limits of type $J$, and finishes proving that $g$ is \texttt{[in]} $\uHom_{!}(\Ind_\pi(C),\mathcal{E})$. Since $g|C\simeq f$, we conclude that the functor (4.4.16.2) is essentially surjective, which finishes proving 4.4.16. Moreover, we have constructed an explicit quasi-inverse functor, $\psi$,
\msp{page 106}
fitting into the commutative diagram
\[
\begin{tikzcd}[column sep=large, row sep=large]
\uHom_{!\,\pi}(C,\mathcal{E}) \arrow[r, "\psi", "(\sim)"'] \arrow[d, hook, "\text{fully faithful incl.}"] & \uHom_{!}(\Ind_\pi(C),\mathcal{E}) \arrow[r] \arrow[d, hook, dotted, "\text{fully faithful}"] & \uHom_{!\,\pi}(\Ind_\pi(C),\mathcal{E}) \\
\uHom(C,\mathcal{E}) \arrow[r, "L", "\substack{\sim\\ \text{(cf. SGA 4 I 8.7.3)}}"'] & \uHom_{!\,\mathrm{filt}}(\Ind(C),\mathcal{E}) \arrow[ur, "R\ \text{(restriction)}"'] &
\end{tikzcd}
\qquad.
\tag*{(4.4.16.6)} % [sic] (4.4.16.6) -> (4.4.16.7)?, the label (4.4.16.6) is already used for Phi on the preceding page
\]

\end{document}
